% !TEX root = ../main.tex
\section{Responsibility Centers -- de 5 centertyper}\label{thr:responsibility-centers}

\paragraph{Definisjon.} Et \emph{responsibility center} er en organisatorisk enhet (avdeling, divisjon, datterselskap) som lederen holdes ansvarlig for via et avgrenset sett av regnskapsmessige måltall. Centertypen bestemmer både hvilke beslutningsrettigheter lederen får (S1) og hvilke prestasjoner som måles og belønnes (S2/S3). BSZ kap. 17 deler virksomhetens enheter i fem typer: \textbf{cost center}, \textbf{expense center}, \textbf{revenue center}, \textbf{profit center} og \textbf{investment center}. [SLIDE]

\subsection*{Forklaring}

Når en virksomhet desentraliserer beslutninger til divisjonsledere, må toppledelsen velge \emph{hvilke variabler lederen kan styre} og \emph{hvordan resultatet skal måles}. Valget av centertype er et samtidig design av (i) beslutningsrettigheter (input-mix, pris, volum, investering), (ii) regnskapsmessig prestasjonsmål (kostnad, omsetning, profit, ROA/EVA) og (iii) belønningsstruktur. Dette er kjernen i det BSZ kaller \emph{organizational architecture}: de tre dimensjonene må passe sammen. [SLIDE/BOK]

Det grunnleggende prinsippet bak inndelingen er \thref{controllability-principle}{Controllability-prinsippet}: en leder bør bare evalueres på utfall hun faktisk kan kontrollere. Hvis en leder ikke har beslutningsmyndighet over for eksempel salgspris, gir det ikke mening å holde henne ansvarlig for omsetning. Sentertypene reflekterer derfor ulik bredde i delegert beslutningsrett. [SLIDE]

\begin{center}
\begin{tabular}{@{}p{2.6cm}p{4.0cm}p{4.0cm}p{3.0cm}@{}}
\toprule
\textbf{Centertype} & \textbf{Lederen kontrollerer} & \textbf{Måles på} & \textbf{Typiske enheter} \\
\midrule
Cost center & Input-mix, produksjonsmetode (gitt output) & Kostnad pr.\ enhet, total kostnad, output pr.\ budsjettkrone & Produksjon, sykehus, jernbane \\
Expense center & Aktivitetsnivå innenfor budsjett & Forbruk vs.\ budsjett (output er vanskelig å måle objektivt) & HR, regnskap, PR, R\&D \\
Revenue center & Salgsmix, salgsinnsats (priser ofte sentralt fastsatt) & Salgsinntekt vs.\ budsjett & Salgsavdelinger, distribusjon \\
Profit center & Input-mix, produktmiks, salgspris & Profit (omsetning $-$ kostnad) & Forretningsdivisjoner med eget P\&L \\
Investment center & Profitcenter + kapitalinvesteringer & ROA, IRR, EVA & Datterselskaper, store divisjoner \\
\bottomrule
\end{tabular}
\end{center}

\subsubsection*{1. Cost center}
Enheten har ansvar for å minimere kostnader gitt et utfallskrav, eller å maksimere output gitt et budsjett. Lederen styrer kombinasjonen av innsatsfaktorer og produksjonsmetode, men har ikke ansvar for omsetning eller pris. Typiske eksempler: produksjonsavdelinger, sykehus, jernbanedrift. [SLIDE]

\emph{Sentralt poeng:} Minimering av gjennomsnittskostnad gir \textbf{ikke} nødvendigvis profitmaksimering -- kostnadsfokus alene kan gi for lav kvalitet eller feil produktmiks. Dette er en gjentakende kritikk av kostnadsstyring i isolasjon. [SLIDE]

\subsubsection*{2. Expense center}
En spesialform av cost center brukt der \emph{output er vanskelig å kvantifisere objektivt}. Lederen får et fast budsjett og evalueres på om aktivitetene leveres innenfor budsjettet, ofte subjektivt. Typiske eksempler: HR, regnskap, PR, R\&D, juridisk avdeling. [SLIDE]

\emph{Sentralt problem:} Når output ikke kan måles, oppstår fare for \emph{empire building} (lederen vil ha større budsjett enn nødvendig), \emph{overconsumption} (penger brukes opp fordi de er der), og generelt vansker med å avgjøre om enheten er produktiv. [SLIDE]

\subsubsection*{3. Revenue center}
Enheten har ansvar for å maksimere omsetning innenfor et gitt budsjett. Lederen styrer salgsinnsats og produktmiks, mens priser og ofte produktportefølje fastsettes sentralt. Typiske eksempler: salgsdivisjoner, distribusjonsavdelinger. [SLIDE]

\emph{Sentralt problem:} Omsetningsmaksimering ($MR = 0$) er ikke det samme som profitmaksimering ($MR = MC$). Et rent revenue center kan derfor presse pris/volum i retninger som er gode for omsetningen, men ikke for selskapets samlede profit -- særlig hvis selgere kan tilby store rabatter eller dyre kundeløsninger. [SLIDE]

\subsubsection*{4. Profit center}
Enheten har ansvar for både kostnader og inntekter. Lederen styrer input-mix, produktmiks og salgspris, og evalueres på regnskapsmessig profit. Typiske eksempler: forretningsdivisjoner med eget resultatregnskap (f.eks.\ en bilprodusents merkevaredivisjon). [SLIDE]

\emph{Forutsetning for at det fungerer:} Lederen må ha \emph{spesifik viden} (lokal kunnskap) som det ville være kostbart å sende oppover, samtidig som divisjonens beslutninger primært påvirker dens egen profit. Når disse forutsetningene brytes, kommer divisjonene i konflikt: \emph{eksternaliteter} mellom divisjoner, \emph{free-riding} på felles merkenavn, og \emph{transfer pricing-konflikter} ved interne leveranser. [SLIDE/BOK]

\subsubsection*{5. Investment center}
Et profitcenter som i tillegg har myndighet over kapitalinvesteringer. Lederen evalueres på avkastning på den investerte kapitalen -- typisk \emph{Return on Assets} (ROA = regnskapsmessig nettoresultat / total kapital i senteret), \emph{IRR}, eller \emph{Economic Value Added} (EVA = NOPAT $-$ WACC $\times$ kapital). [SLIDE]

\emph{Sentralt problem:} Kortsiktige regnskapsmål kan avskrekke gode langsiktige investeringer. Et investment center evaluert kun på årlig ROA kan utsette eller avvise prosjekter med positiv NPV fordi avskrivninger og oppstartskostnader drar ROA ned første år. EVA forsøker å rette dette ved å trekke fra kapitalkostnaden direkte. [SLIDE/BOK]

\subsection*{Kjerneforutsetninger}
\begin{itemize}
\item \textbf{Controllability:} Lederen må faktisk kunne styre de variablene hun evalueres på. Brutt forutsetning $\Rightarrow$ urettferdig og demotiverende.
\item \textbf{Målbarhet:} Det resultatmålet centertypen bygger på må kunne måles tilstrekkelig presist (kostnad, omsetning, profit, kapitalavkastning).
\item \textbf{Spesifik vs.\ generell viden (Jensen \& Meckling):} Desentralisering (profit-/investmentcenter) er attraktivt når lokal kunnskap er kostbar å overføre. Når kunnskapen er sentralt tilgjengelig, peker det mot cost-/expense-center.
\item \textbf{Lav grad av eksternaliteter mellom enheter:} Jo mer divisjonene påvirker hverandre (felles kunder, felles ressurser, interne leveranser), desto mindre egnet er rene profit-/investmentcentre uten transfer pricing-regime.
\item \textbf{Stabilt mål-måleforhold:} Måltallet må gjenspeile reell verdiskaping. Brutt forutsetning gir \emph{gaming} (jf.\ ratchet, horisontproblem).
\end{itemize}

\subsection*{Muligheter (når hver type bør velges)}
\begin{itemize}
\item \textbf{Cost center} -- velges når enheten leverer et standardisert output som er lett å måle, men ikke har ansvar for pris eller marked. Produksjonshaller, IT-drift, lagerterminaler.
\item \textbf{Expense center} -- velges når output er vanskelig å kvantifisere og enheten leverer støtte/innsats snarere enn et målbart produkt. R\&D, HR, regnskap, PR. Budsjettstyring + subjektiv vurdering er det realistiske alternativet.
\item \textbf{Revenue center} -- velges når enheten har innflytelse på salg og kundekontakt, men ikke skal bestemme pris eller produktmiks. Vanlig der toppledelsen vil holde priskontrollen sentralt (merkebeskyttelse, prisdiskriminering).
\item \textbf{Profit center} -- velges når lederen sitter på spesifik lokal viden om både kostnader og marked, og divisjonen kan operere relativt selvstendig. Geografiske divisjoner, produktdivisjoner.
\item \textbf{Investment center} -- velges når også investeringsbeslutninger krever lokal kunnskap (markedsspesifikke prosjekter, ny kapasitet) og enheten er stor nok til at avkastning på kapital kan måles meningsfullt.
\end{itemize}

\subsection*{Begrensninger og kritikk}
\begin{itemize}
\item \textbf{Cost center:} Risiko for at kvalitet, leveringspresisjon og innovasjon ofres for å minimere målte kostnader. Lederen kan velte kostnader over på andre enheter (urealistiske internprisingseffekter).
\item \textbf{Expense center:} Output-måling er subjektiv $\Rightarrow$ rom for politisk spill, empire building og \emph{soft budget constraints}. Vanskelig å avgjøre om enheten er for stor.
\item \textbf{Revenue center:} Maksimering av omsetning $\neq$ profit. Selgere kan gi for store rabatter eller selge til dårlige kunder for å nå omsetningsmål. Trenger ofte komplementære marginmål.
\item \textbf{Profit center:} Skaper interne eksternaliteter (én divisjons valg påvirker en annens profit), \emph{free-riding} på fellesgoder (merkevare, FoU), og strid om internpriser. Risiko for at lokal profit slår fellesnyttig samarbeid i hjel.
\item \textbf{Investment center:} ROA og lignende regnskapsmål er korttidsforvridde -- avskrivninger og engangskostnader straffer langsiktige investeringer. ROA kan også manipuleres ved å selge unna aktiva (\emph{asset stripping}). EVA hjelper, men krever korrekt kapitalkostnad.
\item \textbf{Generelt:} Alle målbaserte centertyper er sårbare for \emph{gaming}, \emph{horisontproblem} og \emph{ratchet-effekt}. Centertypen alene løser ikke insentivproblemet -- den må kombineres med riktig kontraktsdesign (jf.\ informativeness-principle).
\end{itemize}

\subsection*{Modell / oversikt}

Centertypene kan ordnes i en \emph{stige av delegert beslutningsrett}: jo flere kontrollvariabler lederen får, desto bredere måltall trenger evalueringen.

\begin{center}
\begin{tabular}{@{}lcccc@{}}
\toprule
 & \textbf{Kostnader} & \textbf{Inntekter} & \textbf{Profit} & \textbf{Kapital} \\
\midrule
Cost center & \checkmark & & & \\
Expense center & (budsjett) & & & \\
Revenue center & & \checkmark & & \\
Profit center & \checkmark & \checkmark & \checkmark & \\
Investment center & \checkmark & \checkmark & \checkmark & \checkmark \\
\bottomrule
\end{tabular}
\end{center}

\textit{Tabellen viser hvilke størrelser lederen i hver centertype har beslutningsrett og resultatansvar for. [SLIDE -- BSZ kap.\ 17, slide 4]}

\subsection*{Eksempler}
\begin{itemize}
\item \textbf{Equinor (NO):} Letevirksomheten på norsk sokkel kjøres som et \emph{investment center} -- lederne bestemmer letebrønner og felt\-utbygginger og måles på avkastning på investert kapital. Raffinerivirksomheten Mongstad har historisk vært drevet mer som et \emph{cost center}, der lederen styrer effektivitet i prosessen, mens markedspriser settes globalt. [EKS]
\item \textbf{IKEA:} Den enkelte varehussjef driver et \emph{profit center} -- styrer lokal bemanning, kampanjer og produktutvalg innenfor konsernkonsept, og måles på resultatbidrag. Sentralfunksjoner som design (Älmhult) og IT er \emph{expense centers} på fast budsjett. [EKS]
\item \textbf{Maersk:} Shipping-divisjonene Ocean og Logistics \& Services drives som separate profit/investment centers med egne P\&L og kapitalbudsjetter, mens IT- og HR-funksjoner i konsernet er expense centers. [EKS/CASE]
\item \textbf{NSB/Vy (jernbanedrift):} Togdrift har lenge vært et klassisk \emph{cost center} -- billetter og ruter er regulert, så lederne måles på driftskostnad pr.\ togkilometer. Da Vy-konsernet ble organisert i forretningsenheter, ble enkelte selskaper omgjort til profit centers, noe som økte fokuset på kommersielle ruter og endret atferd merkbart. [EKS]
\item \textbf{Salgsdivisjoner i farmasi (f.eks.\ Novo Nordisk):} Lokale salgsteam er gjerne \emph{revenue centers} fordi priser forhandles sentralt og refusjonssystemer er stabile -- men kombineres med marginvilkår for å unngå at selgere skyver lavmargin-produkter. [EKS]
\end{itemize}

\subsection*{Kobling til søyle og andre teorier}
Søyle: \STwo\ (centertypen \emph{er} et prestasjonsmålingsregime), men også \SOne\ (allokering av beslutningsrett) og \SThree\ (måltallet kobles til belønning). Centertyper er derfor et godt eksempel på at de tre søylene må designes samlet.

Relaterte teorier: \thref{controllability-principle}{Controllability-prinsippet} (grunnregelen bak valget), \thref{transfer-pricing}{Transfer pricing} (oppstår nettopp når man har profit-/investmentcentre med interne leveranser), \thref{decision-management-control}{Decision Management vs.\ Decision Control} (centertyper definerer hvem som initierer/implementerer vs.\ ratifiserer/monitorerer), \thref{specific-vs-general-knowledge}{Spesifik vs.\ generell viden} (Jensen \& Meckling-argumentet for når desentralisering lønner seg), \thref{ratchet-effect}{Ratchet-effekten} (gaming-problem som rammer alle målbaserte centre).

\paragraph{Brukt i spørsmål:} Q01, Q02, Q05.
